\documentclass[lettersize,journal]{IEEEtran}
\usepackage{amsmath,amsfonts,amssymb}
\usepackage{algorithmic}
\usepackage{algorithm}
\usepackage{array}
\usepackage[caption=false,font=normalsize,labelfont=sf,textfont=sf]{subfig}
\usepackage{textcomp}
\usepackage{stfloats}
\usepackage{url}
\usepackage{verbatim}
\usepackage{graphicx}
\usepackage{color}
% \usepackage{appendix}
\usepackage{cite}
\newcommand{\parab}[1]{\noindent\textbf{#1} }
\newenvironment{proof}{{\noindent\it Proof:} }{\hfill $\square$\par}
\hyphenation{op-tical net-works semi-conduc-tor IEEE-Xplore} 

\begin{document}
\subsection{Finite-budget operation count}
\label{app:budget-count}
One approximate operation is a scalar Gaussian or finite uniform draw, one
complete three-outcome IRS transition, or one weight coin. Its error includes
the numerical evaluation needed for that draw and all of its bounded random
draws. Guard bits are allocated within an operation so that its conditional
total-variation error is at most $2^{-b}$ uniformly over the input. This is
the premise of~\eqref{eq:cyclic-sampling-error}. A vector draw is counted
coefficient by coefficient. A norm-conditioned Gaussian vector uses at most
$b+3$ trials, as in the bounded sampler above; its cutoff comparison is
included in the same allocation.

A deliberately loose common operation bound is
\begin{align*}
	R_{\rm count}=128\bigl[&4UL_k(Dn+mln+mn+d+4)\\
	&+S((b+3)Dn+2d+8)\\
	&+Q(mln+mn+d+4)\bigr].
\end{align*}
The first term counts base and auxiliary secret trials, public-matrix
draws, local uniform normal-form instances and initialization coins. The
second counts every possible ring or auxiliary branch attempt, its bounded
Gaussian-vector trials, at most $d$ IRS transitions and the weight coin.
The last term overcounts all oracle outputs coefficient by coefficient,
including deterministic programmed entries and all $UK$ bulk installations.
All repeated and preprocessing queries were already included in $Q$.

The factor 128 provides common slack for the bounded sampler-bearing
comparison games and their couplings. There are at most eight such
versions in a displayed subreduction: real signing, capped independent
challenges, private rejected entries erased, direct accepting simulation,
norm-conditioned keys, unconditioned NTWE keys, uniform keys and the
normal-form extraction game. Either candidate hybrid in anonymity and
either primitive branch in unforgeability may each be charged separately;
$8\cdot2\cdot2<128$. Replaying a fork is already included in the aggregate
resource bounds, so it is not an omitted extra execution. The unbounded
analytical continuation used only in the simulation coupling is not run
by a computational reduction, and its cutoff event is charged to
$\delta_{\rm cap}$ instead. The count also covers the simpler bounded
versions that skip rejected signing trials.

For the illustration in Section~VII-B1 of the main paper, exact integer evaluation gives
$R_{\rm count}<2^{75}<R=2^{88}$. For the wider query envelope considered below,
the maximum count is still below $2^{88}$. Thus the same value of $R$ can
be used in both comparisons. 

\subsubsection{Conditional numerical error propagation}
For the finite-budget illustration in Section~VII-B1 of the main paper, the operation count
above with $R=2^{88}$ gives
$\delta_{\rm samp}\le2^{-232},2^{-296},2^{-360}$, conditionally on the
per-operation total-variation premise. The new-tag hash-guessing term
$KQ^2/2^d$ is exactly $2^{-140},2^{-220},2^{-372}$. After propagating all
comparison terms through Theorems 1--4, the sum of their four non-hardness
parts is below $2^{-138},2^{-218},2^{-304}$, respectively. These figures
include the tighter candidate norm-tail and invertibility bounds, the
key-generation and signing cutoffs, oracle simulation, collisions and
finite-precision comparisons.

For example, in this envelope the coefficient of $\sqrt{\epsilon_{\rm M}}$
in the new-tag bound is about $2^{121.03}$, and the NTWE coefficient in the
anonymity bound is $2^{33}$. To make each of these contributions at most a
quarter of $2^{-t}$, sufficient primitive bounds for $t=120,180,256$ are
$\epsilon_{\rm N}\le2^{-155},2^{-215},2^{-291}$ and
$\epsilon_{\rm M}\le2^{-487},2^{-607},2^{-759}$, respectively, at the
augmented reductions' running times. These required advantage bounds use the conditional interpretation in Section~VII-B1 of the main paper.

The dependence on the envelope is material. With aggregate external caps
$q_G,q_R,q_A\le2^{48}$, the same bulk accounting is valid with $K=2^{50}$
and $Q=2^{67}$. The new-tag hash terms then become
$2^{-96},2^{-176},2^{-328}$. Consequently the first two profiles do not
meet their label targets even for the non-hardness part of this more
permissive illustration. These are upper-bound limitations, not attacks.
The accompanying executable budget check records both envelopes and every
affected error and reduction coefficient.
	
\end{document}
